% Ad Familiares 10.15 --- headnote
% ad-familiares-10-15, c. 13 May 43 BC, on the Isère (Allobroges)

\textit{Plancus to Cicero, written from the field on or shortly
after 13 May 43 BC --- Perseus dateline \textit{Scr. in Gallia
Narbonensis cis Isaram circ. iii Id. Mai. a. 711 (43)}. A
dispatch from a moving army. Plancus, governor of Transalpine
Gaul and Gallia Comata, has been edging toward open
intervention against Antony since the news of Mutina; this
letter reports that he has now built a bridge across the
Is\`ere in a single day, crossed his army on 12 May, and sent
his brother Lucius Plancus south on 13 May with four thousand
cavalry to head off Lucius Antonius at Forum Iulii (modern
Fr\'ejus). He himself follows by forced marches with four
light-order legions and the remaining horse.}

\textit{The political news inside the military news is
Lepidus. Lepidus is the wild card of the spring: his province
covers Narbonese Gaul and Nearer Spain, his army is large,
and his loyalties are unread. Plancus claims here, through
the intermediary Marcus Iuventius Laterensis, to have secured
Lepidus's word that if he cannot keep Antony out of his
province, he will fight him --- and that Plancus should come
to join forces. The promise will not hold: within a fortnight
Lepidus's army will go over to Antony and Lepidus himself will
follow. Plancus does not yet know this, but the letter's tone
--- one ruined brigand against the children, the city, and
the state, with Lepidus committed by his word --- is the
optimistic version of a situation that will collapse. The
closing line, \textit{fac valeas meque mutuo diligas}, is the
warm signature of one of the most accomplished political
correspondents of the late republic.}
